% ID: SEL-D0021-OA-E03
% ESTADO: validada. S022, revisada con los criterios oficiales.
Por independencia, $P(A\cap B)=P(A)P(B)$, así que $P(B)=0.25$.
\textbf{a)} $P(A\cap B^c)=0.4-0.1=0.3$. Dividimos por
$P(B^c)=0.75$ y obtenemos $\boxed{P(A\mid B^c)=0.4}$.

\textbf{b)} La incompatibilidad da $P(A\cap C)=0$.
Por De Morgan, $P(A^c\cap C^c)=1-P(A\cup C)=1-0.4-0.2=\boxed{0.4}$.

\textbf{c)} Dentro de $B^c$, el suceso $A^c$ es el complementario
de $A$. Por tanto $P(A^c\mid B^c)=1-0.4=\boxed{0.6}$.
